\section{结论与启示}\label{sec:conclusion}

\subsection{研究结论}
本研究从报告发布、行业收益的实现时间以及训练样本的可用时间出发，检验货币政策文本情绪能否用于 A 股行业收益预测。NTUSD 简体与 15 期归档抽检分值的 Pearson 相关为 $0.695$，与 DLUT 的相关为 $0.829$。每期仅有一份评分，评分过程与独立性仍待核验，评分者间一致性和重测信度均未评估。因此，H1 仅在暂定数值判据下获得初步、部分支持，独立人工效度仍有待验证。

H2 未获支持。主事件窗及不同设定中出现的未校正显著性经多重检验校正后消失。原口径下，主词典模型在共同样本上的样本外 $R^2$ 低于无文本基准。采用严格口径并约束训练时点后，差距仍为负（$\Delta R^2=-0.341$）。H3 暂未获支持。三类行业特征代理在联合截面回归中均未达预设显著性水平，前后子样本行业敏感度排序接近不相关。H4 未获支持。基准提示词下，大语言模型的平均秩相关有所提高，但点预测误差扩大，其联合模型系数也缺乏统计支持。

\subsection{与已有文献的对话}
姜富伟等将货币政策报告文本与市场层面的资产价格联系起来\cite{jiangHuHuang2021}。对全市场共同成分有解释力的季度信号，只有在行业暴露存在稳定差异时才可能改善截面排序。本文对行业敏感度的检验与林雨欣研究央行沟通行业差异的方向相近\cite{linYuxin2025}，联合代理检验与跨期稳定性检验为比较提供了可复核的量化依据。国外研究分别从政策沟通的措辞及货币政策目标推断入手\cite{luccaTrebbi2009,shapiroWilson2022}。沟通载体、制度背景、估计对象和收益频率的差异限制了跨研究比较。本稿的贡献在于以相同报告集比较两类度量方式，并识别收益窗口和可用信息约束对于评判行业预测力的重要性。

\subsection{实践启示、局限与展望}
现有结果不支持将单一文本情绪分数直接用于行业轮动。研究设计可以进一步区分回顾性经济判断与前瞻性政策指引，比较措辞变化、可读性与政策主题等特征，并将跨期验证和收益实现约束写入实验协议。更细的发布时点和交易频率有助于缩短事件窗，独立双人评分及复评可检验报告情绪的测量误差。正则化或树模型需要在相同信息集和共同测试样本上重新评估。

本研究的有效时间单位主要是季度。历史估计窗重叠、部分事件窗中同期政策公告、行业行情起点差异，以及宏观数据未完全实现历史版本控制，均会影响估计或复现。抽检评分的过程记录尚不完整，有待通过独立双人评分和间隔重测补充验证。大语言模型调用的底层修订号亦不可得。以上证据主要用于关联描述与预测评价，适用范围限于现有样本、收益频率和模型。

\subsection{数据与代码可用性}
实验脚本、运行日志、输入清单及图表对应关系见附录\ref{sec:reproduction}。报告原文可由人民银行公开栏目与项目索引定位。完整报告缓存、宏观与利率缓存、外部因子文件和原始模型响应按来源许可与项目记录管理。本文以行业汇总行情和公开政策文本为核心实证材料。署名贡献、资助及利益冲突事项由团队于定稿时核定。人机分工详见赛事单独提交的声明。
